\section{Related Work}
\subsection{Application and Mobile Manipulation}
Vega provides a relevant hazardous-environment precedent: Bird \emph{et al.} describe a small ground robot for nuclear decommissioning with engineering directed toward its operating environment~\cite{bird2022vega}. Its application illustrates the importance of environment-specific qualification. Stretch combines a differential mobile base and compact manipulator, supported by design analysis and physical task demonstrations~\cite{kemp2022stretch}. HomeRobot treats navigation, recognition, grasping and placement as a connected mobile-manipulation problem with distinct simulation and real-world evaluation components~\cite{yenamandra2023homerobot}. These studies motivate HazBot's complete retrieval sequence and the coordination of its sensing and actuation modules.

\subsection{Embedded Integration and Coordination}
ROS 2 provides a common communication and software framework for distributed robotic components~\cite{macenski2022ros2}. Nav2~\cite{macenski2020marathon2} and SLAM Toolbox~\cite{macenski2021slamtoolbox} provide established navigation and mapping capabilities. These components supply reusable interfaces for integrating onboard observations and motion requests. Edge-device detector evaluations further motivate explicit characterization of model and device choices~\cite{alqahtani2026edge}, including the configuration used for local inference.

Command arbitration and sensor-dependent stopping are established in systems such as \code{twist\_mux} and Nav2 Collision Monitor~\cite{twistmux,nav2collision}. HazBot focuses on the connection between a four-subsystem embedded design, retrieval-phase prerequisites and reproducible characterization of the particular geometry and gateways that implement those interfaces. The reference coordinator makes this integration inspectable before a physical mission benchmark is attempted.
